\documentclass[dvipsnames]{article}

\usepackage[utf8]{inputenc}
\usepackage{fontspec}
\usepackage{polyglossia}
\usepackage{unicode-math}
\usepackage{csquotes}
\setdefaultlanguage{russian}
\setotherlanguage{english}

\setmainfont{CMU Serif}
\setsansfont{CMU Sans Serif}
\setmonofont{CMU Typewriter Text}
% \setmathfont{Latin Modern Math}

\usepackage[letterpaper,top=2cm,bottom=2cm,left=3cm,right=3cm,marginparwidth=1.75cm]{geometry}

\usepackage{amsmath}
\usepackage{graphicx}
\usepackage{amsthm}
\usepackage{array}
\usepackage{fancyhdr}
\usepackage{bussproofs}
\usepackage[colorlinks=true, allcolors=blue]{hyperref}
\usepackage[backend=biber]{biblatex}
\usepackage{xcolor}

\usepackage{minted2}
\usemintedstyle{xcode}
\setminted{escapeinside=||}
\newminted[stella]{cpp}{escapeinside=||,frame=single}
\newmintinline[haskell]{haskell}{fontsize=\scriptsize}
\newmintinline[java]{java}{fontsize=\scriptsize}
\newmintinline[python]{python}{fontsize=\scriptsize}
\newmintinline[code]{c}{fontsize=\small}
\newmintinline[smallcode]{c}{fontsize=\small}

\let\smallcodeOrig\smallcode
\renewcommand{\smallcode}[2][]{%
  \ifmmode
    \if\relax\detokenize{#1}\relax
      \mathchoice{\mbox{\normalfont\small\ttfamily #2}}{\mbox{\normalfont\small\ttfamily #2}}{\mbox{\normalfont\small\ttfamily #2}}{\mbox{\normalfont\small\ttfamily #2}}%
    \else
      \mathchoice{\mbox{\normalfont\small\ttfamily #2}}{\mbox{\normalfont\small\ttfamily #2}}{\mbox{\normalfont\small\ttfamily #2}}{\mbox{\normalfont\small\ttfamily #2}}%
    \fi
  \else
    \smallcodeOrig[#1]{#2}%
  \fi}
\let\codeOrig\code
\renewcommand{\code}[2][]{%
  \ifmmode
    \if\relax\detokenize{#1}\relax
      \mathchoice{\mbox{\normalfont\small\ttfamily #2}}{\mbox{\normalfont\small\ttfamily #2}}{\mbox{\normalfont\small\ttfamily #2}}{\mbox{\normalfont\small\ttfamily #2}}%
    \else
      \mathchoice{\mbox{\normalfont\small\ttfamily #2}}{\mbox{\normalfont\small\ttfamily #2}}{\mbox{\normalfont\small\ttfamily #2}}{\mbox{\normalfont\small\ttfamily #2}}%
    \fi
  \else
    \codeOrig[#1]{#2}%
  \fi}

\usepackage{sourcecodepro}
\usepackage{mdframed}

\DeclareMathOperator{\xor}{\mathsf{xor}}
\DeclareMathOperator{\tru}{\mathsf{tru}}
\DeclareMathOperator{\fls}{\mathsf{fls}}

\title{Языки программирования. Семантика и системы типов}
\author{Николай Кудасов}
\date{25 марта 2026}

\addbibresource{../references.bib}

\begin{document}

\pagestyle{fancy}
\fancyhf{}
\chead{Задание выполнил \textcolor{blue}{Павел Григоренко}}
\cfoot{Языки программирования. Семантика и системы типов (весна 2026, Университет ИТМО)}

\section*{Тема 6. Теоретическое задание \textcolor{blue}{(выполнил Павел Григоренко)}}%

\begin{enumerate}
  \item Рассмотрим расширение простого типизированного $\lambda$-исчисления со ссылками (TaPL 13) операцией
    \[
      \code{with}~x = \code{ref}~t_1~\code{do}~t_2,
    \]
    при которой создаётся временная ячейка памяти (с начальным значением равным вычисленному значению $t_1$),
    доступная только на время выполнения блока $t_2$ (и освобождающаяся по завершении блока).
    Покажите, почему даже такая конструкция (без явного оператора \code{free}) ломает безопасность типизации:
    \begin{enumerate}
      \item Выпишите два правила вычисления для конструкции \code{with}--\code{ref} (одно для случая, когда $t_1$ является значением, другое для случая, когда $t_1$ не является значением). Правила должны учитывать состояние памяти ($\mu$).
      \item Предложите правило типизации для конструкции \code{with}--\code{ref}. Правило должно учитывать структуру типизации памяти ($\Sigma$).
      \item Покажите, что сохранение типизации~\cite[Теорема 13.5.3]{TaPL-ru} не выполняется.
        Составьте выражение \code{t}, использующее конструкцию \code{with}--\code{ref}, и представьте состояния памяти $\mu$ и $\mu'$,
        контекст $\Gamma$, структуры типизации $\Sigma$ и $\Sigma'$, а также типы $T$ и $T'$,
        для которых выполняются условия:
        \[
          \code{t} \mid \mu \longrightarrow^* \code{t}' \mid \mu',
        \]
        \[
          \Gamma \mid \Sigma \vdash \mu,
          \qquad
          \Gamma \mid \Sigma \vdash \code{t} : T,
          \qquad
          \Gamma \mid \Sigma' \vdash \mu',
        \]
        но при этом
        \[
          \Gamma \mid \Sigma' \vdash \code{t}' : T', \quad \text{где } T \neq T'.
        \]
    \end{enumerate}

  \color{blue}
  \textbf{Решение.}
  \emph{место для решения}
  \color{black}

  \clearpage

  \item Реализуйте следующие функции над списками~\cite[§11.12]{TaPL-ru}, используя ссылки~\cite[§12]{TaPL-ru},
    и \textbf{не} используя \code{fix} или \code{letrec}~\cite[§11.11]{TaPL-ru}:
    \begin{enumerate}
      \item $\code{length[T]} : \code{List[T]} \to \code{Nat}$
      \item $\code{reverse[A]} : \code{List[A]} \to \code{List[A]}$
      \item $\code{map[A,B]} : (\code{A} \to \code{B}) \to \code{List[A]} \to \code{List[B]}$
    \end{enumerate}
    При необходимости используйте стандартные операции со списками из TaPL (например,
    \code{nil[T]}, \code{cons[T] t u}, \code{isnil[T] t}, \code{head[T] t}, \code{tail[T] t}) и доступные операции со ссылками
    (\code{ref}, \code{!}, \code{:=}, последовательное выполнение).

  \color{blue}
  \textbf{Решение.}
  \begin{enumerate}
      \item \
\begin{minted}[frame=single,linenos,fontsize=\footnotesize]{scheme}
length[T] =
let result = ref (λ l : list[T] . 0) in
(
  f := (λ l : list[T] .
    if isnil[T] l
    then 0
    else succ ((!f) (tail[T] l))
  )
  ;
  ! f
)
\end{minted}
      \item \
\begin{minted}[frame=single,linenos,fontsize=\footnotesize]{scheme}
reverse[T] =
let move = ref (λ to : list[T] . λ from : list[T] . to) in
(
  move := (λ to : list[T] . λ from : list[T] .
    if isnil[T] from
    then to
    else (!move) (cons[T] (head[T] from) to) (tail[T] from)
  )
  ;
  (! move) (nil[T])
)
\end{minted}
      \item \
    \begin{minted}[frame=single,linenos,fontsize=\footnotesize]{scheme}
map[T, U] = λ f : (T -> U) .
let inner = ref (λ l : list[T] . nil[U]) in
(
  inner := (λ l : list[T] .
    if isnil[T] l
    then nil[U]
    else cons[U] (f (head[T] l)) ((!inner) (tail[T] l))
  )
  ;
  ! inner
)
\end{minted}
    \end{enumerate}
  \color{black}

  \clearpage
\end{enumerate}

\section*{Список литературы}
\printbibliography[heading=none]

\end{document}

